\subsection{What the framework delivers: one band, a family of claims}
\label{sec:claims}

The results so far are about scope: when a correction may be used, and when it
may not. A scope result needs a companion showing what the framework yields
when applied, and the available one is a survey-inference finding rather than a
substantive one.

For a country with observed rounds, take the contrast surface
$\theta_{a,b}(t)=F_b(t)-F_a(t)$ over ordered round pairs and core thresholds.
Rounds are independent cross-sections, so each round's sampling units are
resampled within its own strata. One studentised one-sided sup-$t$ critical
value covers the whole surface, so a family of claims of increasing
strength---some adjacent pair declines across the core, the span declines,
every adjacent pair declines---is read off one coverage event, with no
multiplicity spent per claim. Certified claims hold jointly with probability at
least $1-\alpha$, by set inclusion in the band, and the count of claims does not
enter.

Table~\ref{tab:claims} applies this to trust in parliament over European Social
Survey rounds 9--11 at $\alpha=0.10$, in \ClmN{} countries.

\begin{table}[htbp]\centering
\caption{Countries certified, by claim and by how uncertainty is handled. The
first column is the sign of the point estimate and carries no inference; the
second adds design-based uncertainty pointwise; the third adds simultaneity over
the whole contrast surface. Trust in parliament, rounds 9--11,
$\alpha=0.10$.}\label{tab:claims}
\begin{tabular}{lrrr}\toprule
claim & point estimate & pointwise & simultaneous\\\midrule
some adjacent pair declines & \ClmAnyPoint & \ClmAnyPointwise & \ClmAnySim\\
the span declines & \ClmNetPoint & \ClmNetPointwise & \ClmNetSim\\
every adjacent pair declines & \ClmPerPoint & \ClmPerPointwise & \ClmPerSim\\
\bottomrule\end{tabular}
\end{table}

Two reductions compound, and they are worth separating because they answer
different questions. Of the \ClmAnyPoint{} countries whose point estimates fall
in some adjacent pair, \ClmAnyPointwise{} survive design-based uncertainty
applied one contrast at a time, and \ClmAnySim{} survive a band covering the
whole surface at once. Roughly a quarter of the apparent movement is within
sampling error, and a further fifth of the original is within the multiplicity
of looking at every pair and threshold. The claim rung then compounds again:
\ClmAnySim{} countries certify some adjacent decline, \ClmNetSim{} certify the
span (\ClmNetSet), and \ClmPerSim{} certifies decline at every adjacent pair.

We record that the earlier implementation reported the first and third columns
of the first row as a contrast between a ``marginal reading'' and a
simultaneous band. The first column is not a marginal statistical reading; it is
the sign of an estimate with no uncertainty attached. The decomposition above is
the honest form of the same comparison, and it remains a large effect.

\paragraph{What this does not establish.} Nothing here is a claim about
political change. The window contains a pandemic round and a mode transition in
several countries, either of which produces the shifts these bands certify, and
measurement comparability across rounds is assumed rather than shown. The
statement is about response distributions under those assumptions. Its purpose
here is to size what simultaneity and design-based uncertainty cost a repeated
cross-national reading, which is a survey-methodological quantity.
